\section{Trabalhos Relacionados}
\label{trabalhosrelacionados}

Esta seção analisa como a literatura recente comparou algoritmos de agregação sob dados não-IID, com o objetivo de responder a uma pergunta: o que já se sabe sobre as condições em que o FedProx supera o FedAvg, e o que ainda não se sabe? Como este trabalho é uma pesquisa experimental, a análise prioriza trabalhos primários publicados nos últimos cinco anos. Os estudos são organizados primeiro pelo contexto de avaliação (benchmarks com partições controladas, saúde, borda e outros domínios), porque o contexto determina os dados e a forma da heterogeneidade. Depois, são comparados segundo critérios derivados do problema de pesquisa: tipo e nível de heterogeneidade, tratamento de $\mu$, repetição dos experimentos, custo medido e ambiente de execução.

\subsection{Seleção dos estudos e base de evidência}

Os estudos vêm de um levantamento sistemático inspirado no PRISMA \citep{page2021prisma}, com triagem por um único revisor, realizado no IEEE Xplore, no Google Scholar e no OpenAlex (este cobrindo ACM, SpringerLink e ScienceDirect). A busca exigiu os termos de FL, FedProx, FedAvg e não-IID no título ou resumo e um termo de comparação ou avaliação no título. Foram incluídos artigos completos revisados por pares, de 2021 a 2026, com comparação experimental e quantitativa de agregadores sob partições não-IID. Foram excluídos preprints, surveys sem experimentos, propostas de algoritmo novo que usam FedAvg e FedProx apenas como linha de base e aplicações que não analisam a agregação. Dos 50 registros brutos, 18 foram aceitos, e 4 vieram de \textit{snowballing}, totalizando 22 (Tabela~\ref{tab:fluxo}).

\begin{table}[!b]
    \centering\small
    \caption{Fluxo de seleção dos estudos (busca de 12/10/2026).}
    \label{tab:fluxo}
    \begin{tabular}{@{}p{6.6cm}r@{}}
        \toprule
        Etapa & Registros \\
        \midrule
        Retorno bruto (IEEE 16, Scholar 29, Springer 4, ScienceDirect 1, ACM 0) & 50 \\
        Após período (2021--2026) & 49 \\
        Após foco na comparação sob não-IID & 39 \\
        Após exigência de experimentos e métricas & 36 \\
        Após revisão por pares e idioma & 25 \\
        Após remoção de duplicatas & 22 \\
        Incertos, sem resumo acessível & $-4$ \\
        Aceitos na busca & 18 \\
        \textit{Snowballing} & $+4$ \\
        \midrule
        \textbf{Estudos aceitos} & \textbf{22} \\
        \bottomrule
    \end{tabular}
\end{table}

A profundidade da leitura não foi a mesma para todos os estudos, e isso condiciona as conclusões. O NIID-Bench \citep{li2022niidbench}, \citet{quintana2026mammography} e \citet{marella2025fedonco} foram analisados no texto completo. Para a maioria dos demais, de acesso restrito, a análise se apoia no resumo publicado. Por isso, ``não informado'' significa ausência na fonte consultada, e não que o estudo deixou de fazer a medição. Dois estudos aceitos na triagem do levantamento, marcados com $\dagger$, não puderam ser reconferidos em fonte acessível, e por isso não sustentam nenhuma conclusão desta seção. A Tabela~\ref{tab:comparacao} indica a base usada para cada estudo.

\subsection{Apresentação dos trabalhos selecionados}

\paragraph{Benchmarks com partições controladas.}
O NIID-Bench \citep{li2022niidbench} é o trabalho mais próximo do problema investigado. Motivado pela falta de partições representativas nas avaliações anteriores, propõe estratégias para gerar desvio de rótulos, de atributos e de quantidade e compara FedAvg, FedProx, SCAFFOLD e FedNova em nove datasets, com 10 clientes por padrão, 50 rodadas, 10 épocas locais e três repetições, além de um experimento com 100 clientes e fração de participação de 0,1. Conclui que nenhum algoritmo vence em todos os casos, que o caso com uma única classe por cliente é o mais difícil e que, com 100 clientes, o treinamento é muito instável e o SCAFFOLD não funciona bem. \citet{enang2026optimiser} fazem uma comparação controlada de FedAvg, FedProx, FedAdam e DP-FedAvg em CIFAR-10 e MNIST, com partições de Dirichlet e por classe e 20 sementes, e relatam que o FedAdam melhora o desempenho sob heterogeneidade severa, enquanto o FedProx altera a trajetória de treinamento, mas não melhora a estabilidade entre sementes. \citet{ikram2026empirical} analisam FedAvg, FedProx e SCAFFOLD com as mesmas duas famílias de partição, com foco em gargalos práticos, e \citet{mora2022empirical} combinam revisão de soluções e comparação em diferentes níveis de heterogeneidade, em condições de operação iguais; os resumos desses dois estudos não trazem resultados numéricos, o que limita seu uso aqui ao desenho experimental. \citet{sevinc2025distribution} e \citet{ananouch2025optimization} restringem a comparação a FedAvg e FedProx: os primeiros contrastam IID e não-IID em MNIST e Fashion-MNIST, e os segundos variam o número de classes por cliente (1, 2 e 5) com LeNet e MNIST, em 200 rodadas, para emular perfis de IoT, celular e desktop. \citet{onlu2025aggregation} avaliam seis agregadores em cinco datasets visuais, com quatro clientes e \textit{transfer learning}, e \citet{nascimento2024dataskew} avaliam cinco algoritmos em três datasets, com ênfase na classificação facial. Este grupo é o mais próximo do problema investigado, porque controla a heterogeneidade e mede o efeito dos algoritmos em condições iguais.

\paragraph{Aplicações em saúde.}
A motivação neste domínio é a divisão dos dados entre instituições, mas nem sempre os experimentos usam essa divisão real. \citet{quintana2026mammography} usam dados de uma única instituição e criam os clientes por densidade mamária, com dois e quatro clientes, validação cruzada em cinco partes e teste de Wilcoxon, e não encontram benefício estatisticamente significativo de FedProx ou SCAFFOLD sobre o FedAvg. \citet{marella2025fedonco} propõem um benchmark de segmentação tumoral com imagens sintéticas, cinco clientes e três sementes, usando $\mu = 0{,}01$: o FedProx obtém Dice de 0,84 contra 0,85 do FedAvg, diferença não significativa, e menor risco de inferência de pertencimento. Em sentido oposto, \citet{mahmood2025alzheimer} relatam o FedProx como o melhor e o mais robusto no diagnóstico de Alzheimer por ressonância, e \citet{delehouzee2024crosssilo} o apontam como o melhor algoritmo no conjunto para classificação binária de raios X em poucos hospitais. \citet{vajjhula2025har} encontram até 97\% de acurácia para o FedAvg e 91\% para o FedProx no reconhecimento de atividades em cenário não-IID, embora destaquem a estabilidade do FedProx sob privacidade diferencial local. \citet{mim2025healthcare} comparam cinco algoritmos em dados estruturados, imagens e segmentação, simulando saídas de clientes e restrições de banda, e concluem que nenhum domina. \citet{tekrour2026diabetes}$^\dagger$ variam $\mu$ no diagnóstico de diabetes. Para este trabalho, o interesse do grupo está menos nos números, específicos de cada tarefa, e mais em mostrar que os ganhos do FedProx não se repetem de um contexto para outro.

\paragraph{Borda e IoT.}
\citet{carranco2025adhoc} implantam FedAvg, FedProx, FedNova e FedAvgM em uma rede \textit{ad hoc} de Raspberry Pi, com dados tabulares de evasão de clientes: o FedNova se destaca em condições balanceadas, e o FedAvgM, sob distribuições heterogêneas. \citet{hu2025uav} avaliam FedAvg, FedProx, FedOpt e MOON em quatro datasets que representam aplicações de enxames de drones e concluem que a regularização das atualizações locais é importante para a estabilidade sob não-IID. O ParkFL \citep{bakr2026parkfl} compara FedAvg, FedProx, FedOpt e SCAFFOLD na detecção de vagas de estacionamento com dados reais, alcançando mAP de 99,5\% e tempo de comunicação abaixo de 3\% do treinamento. \citet{cheng2026stragglers}$^\dagger$ examinam o efeito de \textit{stragglers} sobre FedAvg e FedProx. Este grupo se relaciona ao trabalho pelo ambiente-alvo, a borda, embora os resumos pouco informem sobre o custo de cada algoritmo no dispositivo.

\paragraph{Outros domínios.}
\citet{ghosh2026renewable} cruzam quatro eixos de heterogeneidade (quantidade, covariáveis, não estacionariedade temporal e disponibilidade) em mais de 190 condições de previsão de energia renovável, com degradação de até 83,4\% no MAE do FedAvg em relação à linha de base centralizada, sob forte desvio de quantidade em dados solares, e apontam a não estacionariedade temporal como o eixo isolado de maior efeito. Em linguagem natural, \citet{julekha2026bangla} relatam superioridade consistente do FedProx sobre FedAvg e FedNova em bengali, enquanto \citet{tan2026instruction}, com cinco clientes e DistilGPT2, obtêm melhor perda e perplexidade com o FedAvg e melhor ROUGE-L com o FedProx.

\subsection{Análise comparativa}

A Tabela~\ref{tab:comparacao} sintetiza as características dos estudos e as compara às deste trabalho. A discussão a seguir interpreta essas diferenças segundo os critérios definidos no início da seção.

\begin{table*}[t]
    \centering\scriptsize
    \newcolumntype{L}[1]{>{\raggedright\arraybackslash}p{#1}}
    \caption{Comparação entre os trabalhos relacionados e a pesquisa proposta. Base: T = texto completo; R = resumo; $\dagger$ = dado da triagem do levantamento, não reconferido em fonte acessível. Heterogeneidade: Dir.\ = Dirichlet; Pat.\ = patológica/por classe; Qtd.\ = quantidade; Atr.\ = atributos. n.i.\ = não informado na fonte consultada. Os resultados não são comparáveis diretamente entre linhas, pois tarefas, dados e métricas diferem.}
    \label{tab:comparacao}
    \setlength{\tabcolsep}{2.5pt}
    \begin{tabular}{@{}L{2.2cm}cL{2.5cm}L{2.1cm}L{0.9cm}L{1.8cm}L{1.2cm}L{1.0cm}L{1.6cm}L{3.0cm}@{}}
        \toprule
        Estudo & Base & Algoritmos & Tarefa / dados & Clientes & Heterogeneidade & $\mu$ & Repet. & Custo / ambiente & Resultado principal \\
        \midrule
        \multicolumn{10}{@{}l}{\textit{Benchmarks com partições controladas}} \\
        \citet{li2022niidbench} & T & FedAvg, FedProx, SCAFFOLD, FedNova & 9 datasets (visão, tabular) & 10; 100 com $C=0{,}1$ & Dir., Pat., Qtd., Atr. & melhor de 4 valores & 3 & tempo/rodada, comunicação & FedProx: em imagens, ganho claro só com 1 classe por cliente; $|\Delta| \leq 0{,}8$ p.p.\ em Dir.(0,5) e Qtd.; maior tempo \\
        \citet{enang2026optimiser} & R & FedAvg, FedProx, FedAdam, DP-FedAvg & CIFAR-10, MNIST & n.i. & Dir., Pat. & n.i. & 20 & n.i. & Não melhora estabilidade entre sementes; FedAdam melhor no severo \\
        \citet{mora2022empirical} & R & n.i. & n.i. & n.i. & Níveis variados & n.i. & n.i. & n.i. & n.i. \\
        \citet{nascimento2024dataskew} & R & 5 algoritmos & 3 datasets, incl. facial & n.i. & Desvios variados & n.i. & n.i. & n.i. & n.i. \\
        \citet{ikram2026empirical} & R & FedAvg, FedProx, SCAFFOLD & 2 datasets & n.i. & Dir., Pat. & n.i. & n.i. & n.i. & n.i. \\
        \citet{sevinc2025distribution} & R & FedAvg, FedProx & MNIST, F-MNIST & n.i. & IID × não-IID & n.i. & n.i. & n.i. & Ligeiramente mais estável \\
        \citet{ananouch2025optimization} & R & FedAvg, FedProx & MNIST, LeNet & n.i. & Pat. (1, 2, 5 classes) & n.i. & n.i. & n.i. & n.i. \\
        \citet{onlu2025aggregation} & R & 6 agregadores & 5 datasets visuais & 4 & n.i. & n.i. & n.i. & n.i. & Consistente; nenhum vence sempre \\
        \midrule
        \multicolumn{10}{@{}l}{\textit{Saúde}} \\
        \citet{quintana2026mammography} & T & FedAvg, FedProx, SCAFFOLD & Mamografia (privado) & 2 e 4 & Densidade mamária & n.i. & 5 partes & n.i. & Sem benefício significativo sobre o FedAvg \\
        \citet{marella2025fedonco} & T & FedAvg, FedProx, FedBN, DP-SGD & TC sintética, segmentação & 5 & Atr. simulados & 0,01 & 3 & n.i. & Dice 0,84 × 0,85; sem diferença significativa \\
        \citet{mahmood2025alzheimer} & R & FedAvg, FedProx, FedSGD, FedBN & Ressonância (Alzheimer) & n.i. & não-IID & n.i. & n.i. & n.i. & Melhor e mais robusto \\
        \citet{delehouzee2024crosssilo} & R & n.i. & Raios X, binária & poucos & Rótulos, Atr., Qtd. & n.i. & n.i. & n.i. & Melhor no conjunto \\
        \citet{vajjhula2025har} & R & FedAvg, FedProx, Krum & UCI HAR & n.i. & IID × não-IID & n.i. & n.i. & n.i. & Inferior ao FedAvg (91\% × 97\%) \\
        \citet{mim2025healthcare} & R & FedAvg, FedProx, SCAFFOLD, FedBN, pFedMe & MIMIC-III, MedMNIST, BraTS & n.i. & não-IID, saídas de clientes & n.i. & n.i. & comunicação & Nenhum domina \\
        \citet{tekrour2026diabetes}$^\dagger$ & $\dagger$ & FedAvg, FedProx, FedSGD & Pima (diabetes) & n.i. & não-IID & 0,1; 0,5; 1 & n.i. & n.i. & n.i. \\
        \midrule
        \multicolumn{10}{@{}l}{\textit{Borda, IoT e outros domínios}} \\
        \citet{carranco2025adhoc} & R & FedAvg, FedProx, FedNova, FedAvgM & Telco Churn & n.i. & IID × não-IID & n.i. & n.i. & Raspberry Pi; custo n.i. & FedAvgM mais robusto \\
        \citet{cheng2026stragglers}$^\dagger$ & $\dagger$ & FedAvg, FedProx & 5 datasets & n.i. & IID a não-IID; \textit{stragglers} & n.i. & n.i. & n.i. & n.i. \\
        \citet{hu2025uav} & R & FedAvg, FedProx, FedOpt, MOON & 4 datasets & n.i. & IID × não-IID & n.i. & n.i. & n.i. & Regularização local é importante \\
        \citet{bakr2026parkfl} & R & FedAvg, FedProx, FedOpt, SCAFFOLD & Vagas (imagens reais) & n.i. & natural & n.i. & n.i. & comunicação & n.i. \\
        \citet{ghosh2026renewable} & R & FedAvg, FedProx, SCAFFOLD & Séries de energia & n.i. & 4 eixos, incl. Qtd. & n.i. & n.i. & n.i. & n.i. \\
        \citet{julekha2026bangla} & R & FedAvg, FedProx, FedNova & PLN (bengali) & n.i. & IID × não-IID & n.i. & n.i. & n.i. & Superior nos dois cenários \\
        \citet{tan2026instruction} & R & FedAvg, FedProx, FedQSN & DistilGPT2, Dolly 15K & 5 & não-IID & n.i. & n.i. & comunicação & Melhor ROUGE-L; perda maior que a do FedAvg \\
        \midrule
        \textbf{Este trabalho} & -- & FedAvg, FedProx, SCAFFOLD, FedAvgM, FedAdam & CIFAR-10, F-MNIST & 100 & Dir.\ (5 $\alpha$), Pat., Qtd. & 0,01; 0,1; 1 (grade exploratória) & 5 (grade), 3 (ref.), 10 (confirm.) & tempo e memória em CPU restrita & Mapa $\alpha \times \mu \times$ (acurácia, custo) \\
        \bottomrule
    \end{tabular}
\end{table*}

\paragraph{Resultados sobre o FedProx.}
Entre os estudos cujas fontes permitem classificar o resultado, não há uma direção única. O FedProx aparece como superior em \citet{mahmood2025alzheimer}, \citet{delehouzee2024crosssilo} e \citet{julekha2026bangla}, e como ligeiramente mais estável em \citet{sevinc2025distribution}. Não supera o FedAvg em \citet{vajjhula2025har}, \citet{quintana2026mammography} e \citet{marella2025fedonco}, e tem resultado misto em \citet{tan2026instruction}. Em \citet{enang2026optimiser}, \citet{carranco2025adhoc} e \citet{mim2025healthcare}, outra classe de algoritmo leva vantagem. Essa contagem indica divergência, mas não mede a eficácia do algoritmo, e os valores de acurácia não podem ser confrontados entre estudos: os 97\% de \citet{vajjhula2025har} e os 0,806 de \citet{mahmood2025alzheimer} vêm de tarefas, dados e modelos sem relação entre si. A comparação válida é a diferença entre algoritmos dentro de cada estudo.

\paragraph{Severidade da heterogeneidade.}
Uma hipótese para parte da divergência, nossa e não dos autores, é o nível de heterogeneidade, mas a evidência é mista. Na Tabela~III do NIID-Bench, nos quatro datasets de imagem, a diferença de acurácia entre FedProx e FedAvg com $\mathrm{Dir}(0{,}5)$ fica entre $-0{,}3$ e $+0{,}5$ ponto percentual, e o FedProx ganha com uma classe por cliente (de $+2{,}3$ a $+17{,}7$), embora nesse caso as acurácias sejam baixas: no CIFAR-10, 12,3\% contra 10,0\%, perto do acaso. A relação, porém, não é monotônica: com duas classes por cliente, a diferença é negativa em MNIST, Fashion-MNIST e SVHN ($-0{,}6$, $-2{,}4$ e $-0{,}9$), e com três classes o sinal varia. Nos dados tabulares, as diferenças mudam de sinal e de magnitude ($+22{,}1$ em rcv1 e $-6{,}3$ em covtype, com $\mathrm{Dir}(0{,}5)$). \citet{quintana2026mammography} sugerem que a heterogeneidade por densidade talvez não seja severa o bastante para justificar métodos específicos, e \citet{vajjhula2025har} obtêm resultados melhores em não-IID do que em IID, o que levanta a dúvida de se a partição produziu o efeito prejudicial pretendido. Por outro lado, \citet{julekha2026bangla} relatam vantagem do FedProx também em IID, o que não se explica por severidade. Portanto, a existência de um limiar de heterogeneidade é uma possibilidade a testar, não um fato estabelecido. Varreduras parciais existem: o NIID-Bench varia o número de classes por cliente, \citet{ananouch2025optimization} fazem o mesmo só com FedAvg e FedProx, e \citet{mora2022empirical} declaram avaliar diferentes níveis. Nas fontes consultadas, contudo, não se encontrou a diferença entre os dois algoritmos ao longo de uma varredura de $\alpha$.

\paragraph{Tipo de heterogeneidade.}
A heterogeneidade não é uniforme entre os estudos. Os benchmarks priorizam partições sintéticas, sobretudo de rótulos, enquanto os estudos em saúde que pudemos ler integralmente dividem os dados por atributo clínico ou de imagem, de forma simulada \citep{quintana2026mammography, marella2025fedonco}. Esses estudos não medem, portanto, o mesmo fenômeno que os benchmarks de rótulos, o que pode explicar parte da discordância. O NIID-Bench observa que os desvios de atributos e de quantidade têm pouca influência na acurácia do FedAvg em classificação. Já \citet{ghosh2026renewable} relatam degradação de até 83,4\% no MAE sob forte desvio de quantidade em dados solares. A aparente contradição não se resolve por comparação direta, pois a tarefa é de previsão e a referência é o treinamento centralizado, mas mostra que o efeito do desvio de quantidade pode depender da tarefa. Em classificação de imagens, a Tabela~III do NIID-Bench aponta para a equivalência ($-0{,}8$ a $+0{,}3$ ponto percentual), embora os autores observem que o FedProx costuma obter a melhor acurácia nesse tipo de desvio, considerando também os dados tabulares; é um caso útil para testar a equivalência de forma estatística.

\paragraph{Tratamento de $\mu$.}
A maioria dos resumos não informa o valor de $\mu$, o que limita a interpretação dos resultados negativos. Entre os textos consultados, \citet{marella2025fedonco} fixam $\mu = 0{,}01$, e o NIID-Bench escolhe o melhor valor entre quatro e relata apenas essa acurácia. Esse procedimento mostra o melhor caso do FedProx, e como $\mu = 0$ equivale ao FedAvg, um $\mu$ muito pequeno tende a aproximar os dois por construção. De fato, os autores observam que o melhor $\mu$ é sempre pequeno e que, por isso, os algoritmos têm convergência e acurácia semelhantes. O que permanece pouco visível é como o desempenho varia com $\mu$ em cada nível de heterogeneidade.

\paragraph{Rigor experimental.}
As repetições variam de três no NIID-Bench e em \citet{marella2025fedonco} a 20 sementes em \citet{enang2026optimiser}. \citet{quintana2026mammography} usam cinco partes de validação cruzada com teste de Wilcoxon e reconhecem que o número pequeno de partes reduz o poder estatístico. Isso importa porque as diferenças discutidas são pequenas: em \citet{marella2025fedonco}, a diferença de 0,01 no Dice é relatada como não significativa, com apenas três sementes, e \citet{enang2026optimiser} tratam a própria estabilidade entre sementes como variável de interesse. Sem repetições e teste adequados, diferenças dessa ordem não sustentam uma classificação entre algoritmos.

\paragraph{Custo e ambiente.}
Os estudos priorizam métricas de qualidade. Nas fontes consultadas, a medida de custo mais comum é a de comunicação \citep{bakr2026parkfl, tan2026instruction, mim2025healthcare}, justamente aquela em que FedAvg e FedProx são iguais por construção; \citet{tan2026instruction} relatam comunicação comparável entre todos os algoritmos avaliados. Entre os 22 estudos, o NIID-Bench é o único caso encontrado com tempo de computação por algoritmo: 73~s contra 133~s por rodada em MNIST e 193~s contra 233~s em CIFAR-10, sem descrição do hardware usado. \citet{carranco2025adhoc} executam os algoritmos em Raspberry Pi, mas o resumo não relata o custo por algoritmo. Os dois trabalhos tocam em partes da pergunta sobre custo, sem respondê-la: um mede o tempo sem isolar o termo proximal, e o outro usa hardware restrito sem relatar, nas fontes consultadas, essa medição.

\subsection{Principais contribuições dos estudos existentes}

Individualmente, o NIID-Bench oferece as estratégias de particionamento e a evidência de que o tipo de desvio altera o resultado; \citet{enang2026optimiser} mostram a viabilidade de uma comparação com muitas sementes e a vantagem do otimizador adaptativo sob heterogeneidade severa; \citet{quintana2026mammography} oferecem um exemplo de análise estatística pareada e um resultado negativo bem documentado; \citet{carranco2025adhoc} demonstram a implantação em dispositivos de baixo custo; \citet{ghosh2026renewable} propõem a análise por eixos de heterogeneidade; e \citet{marella2025fedonco} e o ParkFL \citep{bakr2026parkfl} disponibilizam plataformas reprodutíveis.

No conjunto, três resultados se sustentam em mais de um estudo. Primeiro, nenhum algoritmo domina em todos os cenários \citep{li2022niidbench, onlu2025aggregation, mim2025healthcare}. Segundo, o treinamento federado degrada sob não-IID, mas o tamanho da degradação depende da tarefa e da forma do desvio \citep{li2022niidbench, sevinc2025distribution, hu2025uav, ghosh2026renewable}. Terceiro, em nossa interpretação, o FedAvg é um ponto de comparação mais difícil de superar do que a motivação dos algoritmos especializados sugere: pode igualar alternativas mais complexas em alguns contextos \citep{li2022niidbench, quintana2026mammography, marella2025fedonco}, embora outros estudos relatem o contrário \citep{mahmood2025alzheimer, delehouzee2024crosssilo, julekha2026bangla}.

\subsection{Limitações dos trabalhos e oportunidades de investigação}

\paragraph{Limitações declaradas pelos autores.}
\citet{quintana2026mammography} reconhecem o uso de apenas cinco partes, por restrição de recursos, e indicam como trabalho futuro redes federadas maiores e o estudo combinado de fontes de heterogeneidade, como densidade, tamanho do conjunto e distribuição de lesões. \citet{marella2025fedonco} declaram o uso de fatias 2D sintéticas, que não reproduzem toda a complexidade de exames reais, e a participação de todos os clientes em cada rodada, ao contrário de sistemas reais. O NIID-Bench considera a estabilidade do treinamento sob dados não-IID e o cenário de larga escala com poucos dados por cliente problemas em aberto, e relata instabilidade sob participação parcial. \citet{sevinc2025distribution} apontam a necessidade de simulações que reflitam melhor as condições reais.

\paragraph{Limitações identificadas nesta análise.}
Três limitações são interpretações nossas. A primeira é a escala: vários estudos usam poucos clientes (quatro em \citet{onlu2025aggregation}, cinco em \citet{tan2026instruction} e \citet{marella2025fedonco}, dois e quatro em \citet{quintana2026mammography}, e dez por padrão na maior parte dos experimentos do NIID-Bench), o que aproxima seus resultados do cenário \textit{cross-silo} e limita a generalização para a borda. A segunda é a seleção de $\mu$ pela melhor acurácia, que impede avaliar a sensibilidade ao hiperparâmetro. A terceira é o custo: o tempo por rodada do NIID-Bench combina todo o treinamento local e depende da implementação, sem separar o termo proximal e sem considerar quantas rodadas o algoritmo economiza. Além disso, quase toda a evidência vem de simulação, e \citet{valli2026beyond}, estudo metodológico que não integra os 22 por não atender aos critérios de seleção, mostram que simulações podem levar a classificações de configurações diferentes das obtidas em implantações reais.

\paragraph{Da limitação à lacuna.}
Uma limitação isolada não caracteriza uma lacuna. Considerando o conjunto, e com a ressalva de que a maioria dos estudos foi consultada pelo resumo, três aspectos parecem insuficientemente investigados: (i) como a \emph{diferença} entre FedProx e FedAvg varia ao longo de $\alpha$ e de $\mu$, de forma conjunta, já que as varreduras encontradas são parciais; (ii) se o custo local do FedProx é compensado pela redução de rodadas, o que exige medir os dois na mesma avaliação; e (iii) se o efeito do FedProx muda quando rótulos e quantidade são separados, com poder estatístico para afirmar também a ausência de diferença. Não se afirma que esses aspectos nunca foram abordados, mas que não foram encontrados de forma combinada nos estudos analisados. A leitura integral dos estudos de acesso restrito pode alterar esse quadro, e por isso continuará durante o trabalho.

\subsection{Posicionamento da pesquisa proposta}

Os trabalhos mais próximos são o NIID-Bench, pelos tipos de desvio, pela medição de tempo, pelo uso de CIFAR-10 e por seu experimento com 100 clientes e fração de 0,1, a mesma configuração adotada aqui, em que relata forte instabilidade; \citet{enang2026optimiser}, pelos dados, pelas partições e pelo uso de muitas sementes; e \citet{sevinc2025distribution} e \citet{ananouch2025optimization}, por focarem em FedAvg e FedProx. Este trabalho se distingue deles por concentrar o desenho na pergunta sobre quando o FedProx compensa, e não em uma classificação geral de algoritmos.

A pesquisa \emph{avalia e estende} abordagens existentes, sem propor um algoritmo novo. Dela se aproveitam as estratégias de particionamento do NIID-Bench, a repetição com múltiplas sementes de \citet{enang2026optimiser} e o teste pareado de Wilcoxon, como em \citet{quintana2026mammography}. O avanço pretendido está em quatro pontos: relatar todos os valores de $\mu$ para cada $\alpha$, em vez de apenas o melhor; usar teste de equivalência, para que um resultado de ausência de diferença tenha sustentação estatística; incluir um cenário só com desvio de quantidade, para separá-lo do desvio de rótulos; e medir, em CPU restrita, o custo por época e o custo total até a acurácia-alvo. Esse avanço tem limites conhecidos. A grade completa $\alpha \times \mu$ usa cinco sementes, número que não permite significância no teste de Wilcoxon, e por isso será analisada de forma descritiva. A análise confirmatória, com dez sementes, usa um único $\mu^*$, escolhido pela melhor acurácia na fase exploratória; é o mesmo tipo de seleção criticado acima, mitigado por usar sementes novas. Os algoritmos de referência têm só três sementes e servem apenas para situar a ordem de grandeza dos efeitos.

A análise também refinou a proposta em três pontos. Primeiro, o objetivo passa a ser verificar \emph{se} existe um limiar de heterogeneidade e, se existir, localizá-lo. Como o melhor $\mu$ no NIID-Bench foi sempre pequeno e as diferenças em imagens com $\mathrm{Dir}(0{,}5)$ ficaram abaixo de 1 ponto percentual, $\alpha = 0{,}5$ passa a ser também um caso candidato à equivalência. Segundo, a expectativa de superioridade do FedProx na partição patológica com duas classes por cliente não encontra apoio nos resultados do NIID-Bench, e o experimento a trata como uma previsão a ser testada, que pode falhar. Terceiro, o \textit{overhead} medido no NIID-Bench obriga a tratar como hipótese a ser testada, e não como premissa, a ideia de que o custo do termo proximal é pequeno. A pesquisa mantém limitações próprias: trata de classificação de imagens em simulação, com CPU restrita como aproximação de um dispositivo de borda, e seus resultados não se generalizam automaticamente para os domínios discutidos acima.

\subsection{Síntese da análise comparativa}

A literatura recente compara algoritmos de agregação por duas abordagens principais: benchmarks com partições controladas, que permitem isolar fatores, e estudos de aplicação, mais próximos de um domínio real, mas específicos dele e nem sempre baseados em divisões reais entre instituições. As duas convergem ao indicar que nenhum algoritmo domina, e divergem quanto ao FedProx em uma medida que não se explica por uma única causa: diferem as tarefas, a severidade da heterogeneidade, o valor de $\mu$ e o rigor estatístico.

Em classificação de imagens, a evidência mais detalhada, a do NIID-Bench, sugere um quadro mais estreito do que o sugerido pela motivação original do algoritmo: ganho quando cada cliente tem uma única classe, ainda que com acurácias baixas, diferenças pequenas e de sinal variável nos demais casos, proximidade com o FedAvg quando $\mu$ é pequeno e custo de computação potencialmente relevante. É um único estudo, com três repetições, e esse quadro não foi encontrado mapeado ao longo de $\alpha$ nem confrontado com o custo total no cliente. É nessa lacuna que a pesquisa se posiciona: o mapa $\alpha \times \mu \times$ (acurácia, custo) pretende transformar observações dispersas em recomendação de uso, indicando se há condições em que vale adotar o FedProx em dispositivos de borda e em quais o FedAvg, ou uma correção no servidor, é suficiente.
